
\medskip
\phantomsection\label{proof:sga3-V-4.1-b}%
\noindent\textbf{b)} Consider two prime ideals \(x\) and \(y\) of
\(A_0\). We shall show that the equality
\[
  x\cap B=y\cap B
\]
implies the existence of a prime ideal \(z\) of \(A_1\) such that
\[
  x=d_0(z),\qquad y=d_1(z).
\]

Indeed, if this assertion were false, \(x\) would be distinct from
\(\delta_0^{-1}(t)\) for every prime ideal \(t\) of \(A_1\) such that
\(\delta_1^{-1}(t)=y\). For every such \(t\), one would have
\[
  \delta_0^{-1}(t)\cap B
  =\delta_1^{-1}(t)\cap B
  =y\cap B=x\cap B.
\]
It would follow from Cohen--Seidenberg (see Bourbaki,
\emph{Alg\`ebre commutative}, V, \S2, Corollary~1 to Proposition~1)
that \(x\) is contained in none of the ideals
\(\delta_0^{-1}(t)\).%
\footnote{\label{note:sga3-V-21}Editor's note: the original has been expanded in what follows;
in particular, Lemma~\hyperref[lem:sga3-V-4.1.1]{4.1.1}, taken from \hyperref[bib:sga3-V-DG70]{[DG70]}, III,
\S2.4, Lemma~4.3, has been added.}
There are at most \(n\) prime ideals \(t\) of \(A_1\) satisfying
\(\delta_1^{-1}(t)=y\) (see \emph{loc. cit.}, Proposition~3). Hence,
by the Prime Avoidance Lemma (\emph{loc. cit.}, II, \S1,
Proposition~3), there would be an element \(a\in x\) belonging to none
of the ideals \(\delta_0^{-1}(t)\).

Consequently, \(\delta_0(a)\) would belong to none of these prime
ideals \(t\), and \hyperref[lem:sga3-V-4.1.1]{the lemma below} would imply that the norm
\(N_{\delta_1}(\delta_0(a))\) does not belong to \(y\). Here the norm
is computed by regarding \(A_1\) as an \(A_0\)-algebra through
\(\delta_1\); with the notation of \hyperref[proof:sga3-V-4.1-a]{part~a)},
\[
  N_{\delta_1}(\delta_0(a))=\sigma_n.
\]
On the other hand,
\[
  (-1)^{n-1}\sigma_n
  =a^n+\sum_{i=1}^{n-1}(-1)^i\sigma_i a^{n-i},
\]
so this norm belongs to \(B\cap x=B\cap y\), a contradiction.

\medskip
\phantomsection\label{lem:sga3-V-4.1.1}%
\noindent\textbf{Lemma 4.1.1.}---
\textit{Let \(A\to A'\) be a homomorphism of commutative rings which
makes \(A'\) a projective \(A\)-module of rank \(n\). Let
\(p\in\operatorname{Spec}(A)\), let \(q_1,\ldots,q_r\) be the points of
\(\operatorname{Spec}(A')\) lying above \(p\), and let \(a\in A'\).
Then \(a\in q_1\cup\cdots\cup q_r\) if and only if its norm \(N(a)\)
belongs to \(p\).}

Indeed, after replacing \(A\) and \(A'\) by the localizations \(A_p\)
and \(A'_p\), we reduce to the case in which \((A,p)\) is local and
\(A'\) is semilocal, with
\[
  \operatorname{Spec}(A')=\{q_1,\ldots,q_r\}.
\]
In this case \(A'\) is a free \(A\)-module of rank \(n\) (see
Bourbaki, \emph{Alg\`ebre commutative}, II, \S3.2, Corollary~2 to
Proposition~5), and \(N(a)\) is the determinant of the endomorphism
\[
  \ell_a:A'\longrightarrow A',\qquad a'\longmapsto aa'.
\]
Thus one has the equivalences
\[
  N(a)\notin p
  \quad\Longleftrightarrow\quad N(a)\text{ is invertible}
  \quad\Longleftrightarrow\quad \ell_a\text{ is invertible}
  \quad\Longleftrightarrow\quad
  a\notin q_1\cup\cdots\cup q_r.
\]
